\documentclass{article}
\usepackage[T1]{fontenc}
\usepackage{lmodern}
\usepackage{graphicx}
\usepackage{amsmath,amssymb}
\usepackage[a4paper,margin=2cm]{geometry}
\usepackage[absolute,overlay]{textpos}
\setlength{\TPHorizModule}{1mm}
\setlength{\TPVertModule}{1mm}
\usepackage[indonesian]{babel}
\usepackage{hyperref}
\setlength{\emergencystretch}{2em}
% unit: Halaman 15

\begin{document}

% === Halaman 15 (python/native) ===

2.  Serangan berbasis statistik (statistical attack)

•

Menggunakan analisis matematik pada citra untuk menemukan 

perbedaan antara cover image dengan stego image. 

•

Didasarkan pada fakta bahwa penyembunyian pesan ke dalam media 

menimbulkan artefak yang dapat dideteksi secara statistik sehingga dapat 

mengungkap penyembunyian pesan atau pesan yang disembunyikan itu 

sendiri.

•

Contoh serangan statistik: 


a. histogram analysis


b. Regular-singular (RS) analysis


c. Chi-square analysis 


 d. Sample pair (SP) analysis 

15

\end{document}
